\section{Программная реализация и экспериментальное исследование}

\subsection{Программная реализация библиотеки}

\subsubsection{Структура проекта и организация исходного кода}

Проект организован в соответствии с общепринятыми соглашениями
о структуре библиотек на языке C++. Корневой
каталог содержит четыре подкаталога, каждый из которых несёт
единственную ответственность. Файловая структура проекта
представлена в листинге~\ref{lst:project-structure}.

\begin{lstlisting}[caption={Файловая структура проекта},
                   label={lst:project-structure}]
implementation
├── benchmarks
│   └── bench_thread_pools.cpp
├── CMakeLists.txt
├── include
│   └── thread_pool
│       ├── simple_thread_pool.hpp
│       ├── work_stealing_deque.hpp
│       └── work_stealing_pool.hpp
├── src
│   ├── simple_thread_pool.cpp
│   └── work_stealing_pool.cpp
└── tests
    ├── test_simple_pool.cpp
    ├── test_work_stealing_deque.cpp
    └── test_work_stealing_pool.cpp
\end{lstlisting}

Каталог \texttt{include/thread\_pool/} содержит заголовочные
файлы всех публичных компонентов библиотеки. Размещение
заголовочных файлов в подкаталоге \texttt{thread\_pool/}
обеспечивает пространство имён на уровне файловой системы:
пользователь подключает компоненты директивой
\texttt{\#include "thread\_pool/simple\_thread\_pool.hpp"},
что исключает конфликты имён с заголовочными файлами
других библиотек.

Каталог \texttt{src/} содержит файлы реализации. Разделение
интерфейса (заголовочные файлы) и реализации (файлы
\texttt{.cpp}) является стандартной практикой в C++: оно
сокращает время компиляции за счёт раздельной компиляции
модулей и скрывает детали реализации от пользователя
библиотеки.

Каталог \texttt{tests/} содержит автоматические тесты
корректности, построенные на основе библиотеки Google Test.
Каждый модуль библиотеки имеет соответствующий файл тестов:
\texttt{test\_simple\_pool.cpp} тестирует \texttt{SimpleThreadPool},
\texttt{test\_work\_stealing\_deque.cpp} тестирует
\texttt{WorkStealingDeque}, \texttt{test\_work\_stealing\_pool.cpp}
тестирует \texttt{WorkStealingPool}.

Каталог \texttt{benchmarks/} содержит тесты производительности,
построенные на основе библиотеки Google Benchmark. Все сценарии
нагрузки сосредоточены в единственном файле
\texttt{bench\_thread\_pools.cpp}, что упрощает сравнение
результатов двух реализаций.

Система сборки проекта построена на основе CMake версии
3.14 и выше. Файл \texttt{CMakeLists.txt} решает четыре
задачи: определяет библиотеку \texttt{thread\_pool} как
цель сборки, подключает сторонние зависимости посредством
механизма \texttt{FetchContent}, определяет исполняемые
файлы тестов и бенчмарков и настраивает параметры компиляции.

Механизм \texttt{FetchContent} автоматически загружает
и собирает библиотеки Google Test и Google Benchmark
при первой сборке проекта, что исключает необходимость
их ручной установки. Ключевые фрагменты файла
\texttt{CMakeLists.txt} представлены в листинге~\ref{lst:cmake}.

\begin{lstlisting}[language=C++,
                   caption={Ключевые фрагменты CMakeLists.txt},
                   label={lst:cmake}]
cmake_minimum_required(VERSION 3.14)
project(thread_pool VERSION 1.0.0 LANGUAGES CXX)

set(CMAKE_CXX_STANDARD 17)
set(CMAKE_CXX_STANDARD_REQUIRED ON)
set(CMAKE_CXX_EXTENSIONS OFF)

# Основная библиотека — без внешних зависимостей
add_library(thread_pool
    src/simple_thread_pool.cpp
    src/work_stealing_pool.cpp
)
target_include_directories(thread_pool PUBLIC include)

# Автоматическая загрузка зависимостей
include(FetchContent)
FetchContent_Declare(googletest ...)
FetchContent_Declare(googlebenchmark ...)
FetchContent_MakeAvailable(googletest googlebenchmark)

# Тесты — с ThreadSanitizer для обнаружения гонок данных
add_executable(tests ...)
target_compile_options(tests PRIVATE -fsanitize=thread)
target_link_options(tests PRIVATE -fsanitize=thread)

# Бенчмарки — с оптимизацией O3
add_executable(benchmarks ...)
target_compile_options(benchmarks PRIVATE -O3)
\end{lstlisting}

Тесты корректности собираются с включённым Thread Sanitizer
(\texttt{-fsanitize=thread}) --- инструментом динамического
анализа, обнаруживающим гонки данных во время выполнения
программы. Бенчмарки производительности собираются с уровнем
оптимизации \texttt{-O2} и без санитайзеров, поскольку
санитайзеры вносят значительные накладные расходы и искажают
результаты измерений.

Сборка и запуск проекта осуществляются следующей
последовательностью команд, представленной
в листинге~\ref{lst:build}.

\begin{lstlisting}[caption={Сборка и запуск проекта},
                   label={lst:build}]
# Конфигурация в режиме Release
cmake -B build -DCMAKE_BUILD_TYPE=Release

# Сборка всех целей
cmake --build build

# Запуск тестов корректности
ctest --test-dir build --output-on-failure

# Запуск бенчмарков с выводом в JSON
./build/benchmarks \
    --benchmark_format=json \
    --benchmark_out=results.json
\end{lstlisting}

\subsubsection{Особенности программной реализации}

В ходе практической реализации библиотеки был выявлен ряд
нетривиальных проблем, решение которых потребовало отклонения
от первоначального проектного решения. Данный раздел описывает
наиболее значимые из них.

Первая проблема --- гонка данных в деструкторе
\texttt{SimpleThreadPool}. В исходной реализации деструктор
устанавливал флаг остановки и вызывал \texttt{notify\_all()}
без захвата мьютекса:

\begin{lstlisting}[language=C++,
                   caption={Исходная некорректная реализация
                            деструктора},
                   label={lst:destructor-bug}]
// Некорректно: возможна потеря уведомления
SimpleThreadPool::~SimpleThreadPool() {
    stop_.store(true);      // флаг установлен
    cv_.notify_all();       // уведомление отправлено
    for (auto& w : workers_) w.join();
}
\end{lstlisting}

При данной реализации возникала следующая гонка: рабочий поток
мог проверить предикат переменной условия, не обнаружить задач,
затем деструктор устанавливал флаг и отправлял уведомление,
и лишь после этого поток входил в \texttt{cv\_.wait()}.
В результате уведомление было потеряно и поток блокировался
навсегда. Проблема была обнаружена при стресс-тестировании:
тест \texttt{ThreadCountIsCorrect} периодически зависал при
запуске совместно с другими тестами.

Решение состоит в захвате мьютекса перед установкой флага
остановки, что гарантирует атомарность операции с точки зрения
рабочих потоков:

\begin{lstlisting}[language=C++,
                   caption={Исправленная реализация деструктора},
                   label={lst:destructor-fix}]
SimpleThreadPool::~SimpleThreadPool() {
    {
        // Захват мьютекса гарантирует что поток
        // либо уже в cv_.wait(), либо увидит флаг
        // при следующем захвате мьютекса
        std::lock_guard<std::mutex> lock(queue_mutex_);
        stop_.store(true);
    }
    cv_.notify_all();
    for (auto& w : workers_) w.join();
}
\end{lstlisting}

Вторая проблема --- переполнение беззнакового типа в методе
\texttt{pop()} класса \texttt{WorkStealingDeque}. В исходной
реализации алгоритма Чейза-Лева метод \texttt{pop()} первым
действием уменьшал значение \texttt{bottom\_} на единицу:

\begin{lstlisting}[language=C++,
                   caption={Ошибка переполнения в методе pop()},
                   label={lst:pop-bug}]
bool pop(std::function<void()>& task) {
    // Ошибка: если bottom_ == 0, то bottom_ - 1
    // равно 18446744073709551615 (переполнение size_t)
    std::size_t bottom =
        bottom_.load(std::memory_order_relaxed) - 1;
    ...
}
\end{lstlisting}

При вызове \texttt{pop()} на пустой очереди, когда
\texttt{bottom\_\,=\,0}, беззнаковое вычитание приводило
к переполнению и обращению к недопустимой области памяти,
что вызывало аварийное завершение программы (segmentation fault).
Проблема была обнаружена тестом \texttt{PopOnEmptyReturnsFalse}.

Решение состоит в проверке равенства \texttt{bottom\_} и
\texttt{top\_} до выполнения каких-либо изменений:

\begin{lstlisting}[language=C++,
                   caption={Исправленный метод pop()},
                   label={lst:pop-fix}]
bool pop(std::function<void()>& task) {
    std::size_t bottom = bottom_.load(std::memory_order_relaxed);
    std::size_t top    = top_.load(std::memory_order_acquire);

    // Проверяем до вычитания — исключаем переполнение
    if (bottom == top) {
        return false;
    }

    bottom -= 1;
    ...
}
\end{lstlisting}

Третья проблема --- отказ от lock-free реализации
\texttt{WorkStealingDeque}. Первоначально планировалась
реализация неблокирующей двусторонней очереди по алгоритму
Чейза-Лева~\cite{Chase2005} на основе атомарных операций
\texttt{std::atomic}. Однако в ходе тестирования были обнаружены
ошибки типа \texttt{double free} и повреждение кучи
(\texttt{malloc\_consolidate(): invalid chunk size}),
связанные с гонкой данных при освобождении памяти под задачи.

Корень проблемы состоит в том что вор читает указатель на задачу
из буфера до выполнения операции CAS. В промежутке между чтением
и CAS владелец может выполнить \texttt{pop()} и удалить задачу,
оставив у вора висячий указатель. Корректное решение требует
применения механизмов безопасного освобождения памяти в
конкурентной среде --- hazard pointers или эпохального удаления
(epoch-based reclamation), реализация которых существенно
выходит за рамки настоящей работы.

В качестве рабочего варианта была выбрана реализация на основе
\texttt{std::deque} и \texttt{std::mutex}, корректность которой
очевидна и легко верифицируется. Данное решение не влияет
на корректность сравнительного исследования, поскольку обе
реализации пулов потоков используют одну и ту же структуру
\texttt{WorkStealingDeque}.

Для обнаружения гонок данных в ходе разработки использовался
инструмент Thread Sanitizer (TSan) --- динамический анализатор,
встроенный в компиляторы GCC и Clang. TSan инструментирует
код во время компиляции и обнаруживает гонки данных во время
выполнения программы. Тесты корректности собираются с флагами
\texttt{-fsanitize=thread}, что позволяет автоматически
проверять отсутствие гонок данных при каждом запуске тестов.
В финальной версии библиотеки Thread Sanitizer не обнаружил
ни одной гонки данных.